% !TeX root = ../main.tex

\section{Metodología}

\subsection{El método de Van Noort}
El método original planteado por Van Noort plantea lo siguiente. Imagina que tenemos unos datos observados $d$ y un modelo teórico $F(x)$ que depende de unos parametros $x$, que en nuestro caso ese modelo sería la aproximación de campo débil y el parámetro el campo magnético paralelo. SI partimos de una estimación inicial de los parámetros $x_0$ el error sería:

\begin{equation}
    \label{eq:beta}
    \beta = d -F(x_0)
\end{equation}

Para mejorar nuestra estimación queremos dar una correción $\nu$, de modo que nuestra nueva estimación sea $x_0 + \nu$

Dado que, según el paper, $F(x)$ suele ser una función dificil y no lineal, se aproxima a una serie de Taylor:

\begin{equation}
    F(x_o + \nu) \simeq F(x_0) + J \nu
\end{equation}

Donde $J$ es la matriz Jacobiana del sistema. \\

El objetivo ideal es que la nueva estimación encaje perfectamente con los datos observados, siendo el error cero:

\begin{equation}
    \begin{split}
        d - F(x_o + \nu) &= 0\\
        d - (F(x_0) + J \nu) &= 0\\
        J \nu = \beta
    \end{split}
\end{equation}

Donde $\beta = d - F(x_0)$. El problema con los datos experimentales que tenemos muchos mas puntos que parametros a ajustar. Eso significa que $J$ es una matriz rectangular. Esto significa que es una matriz sobredeterminada. \\

Para resolverlo usamos el enfoque de minimos cuadrados. Tal como lo proponen en el paper, al realizar este enfoque nos sale el siguiente sistema de ecuaciones:

\begin{equation}
    \label{eq:sistemaResolver}
    J^T J \nu = J^T \beta
\end{equation}

Gracias a esto podemos resolver esta matriz porque $J^T J$ si que es cuadrada y resoluble.

\subsection{Adaptarlo a nuestro caso}
La matriz $J$ es la matriz jacobiana del sistema físico que estamos representando. En nuestro caso implicaría dos transformaciones. Tanto el cálculo del campo magnético a partir de la aproximación débil cómo la convolución con la PSF. Eso implicaría que cada parámetro de $x_i$ sería el enlace pixel a pixel de la imagen original. Si tenemos una imagen de 1400x1400 entonces tendríamos una matriz Jacobiana tan grande que ocuparía terabytes de memoria ram.\\

Aquí es cuando podremos aprovechar nuestro sistema. Si definimos nuestro J para nuestro sistema como un operador en vez de una matriz y descubrimos que es lineal, entonces el sistema estaría perfectamente definido y podríamos resolver la ecuación (\ref{eq:sistemaResolver}).